%===============================================================================
% SEÇÃO 4 — Descrição do Processo e da Planta Piloto
%===============================================================================

\section{Descrição do Processo e da Planta Piloto}
\label{sec:processo}

\subsection{Configuração da planta piloto}

A planta piloto analisada neste trabalho está instalada no LabMater/UFPR e é
dedicada à reforma a seco de \emph{biogás}, utilizando um reator de leito fixo
catalítico aquecido por forno elétrico.\citep{oliveira2021,labmater2024}
O reator consiste em um tubo metálico inserido em um forno de aquecimento
resistivo, carregado com um leito de catalisador à base de níquel
(por exemplo, 20Ni/Si-MCM-41), dimensionado para operar com vazões da ordem de
dezenas de litros por hora de biogás sintético sob pressão
atmosférica.\citep{oliveira2021} Em condições nominais, a planta opera entre
650 e 800~\(^\circ\)C, com conversão global de biogás superior a 97~\% e
produção de cerca de 30~L\(\cdot\)h\(^{-1}\) de hidrogênio em regime
contínuo.\citep{labmater2024}

\begin{figure}[ht]
\begin{center}
% === INSERIR FIGURA: Diagrama P&ID simplificado ===
% \includegraphics[width=8.4cm]{diagrama_pid_planta.eps}
\fbox{\parbox{7.5cm}{\centering\vspace{1.2cm}
{\small [FIGURA A INSERIR]\\[4pt]
Diagrama P\&ID simplificado da planta piloto}\vspace{1.2cm}}}
\caption{Diagrama P\&ID simplificado da planta piloto de BDR, LabMater/UFPR.}
\label{fig:planta}
\end{center}
\end{figure}

\subsection{Instrumentação e arquitetura de controle}

A planta piloto é composta por três unidades de aquecimento elétrico em série:
um forno pré-aquecedor e dois fornos reatores. O forno pré-aquecedor é governado
por uma única malha de controle de temperatura. Cada forno reator é operado por
três malhas independentes, totalizando sete malhas de controle automático
distribuídas pela planta.

O sistema de controle é implementado em um Controlador Lógico Programável (CLP)
Siemens S7-1200, que centraliza a aquisição dos sinais de instrumentação
analógica, executa os algoritmos de controle e comanda os atuadores de
aquecimento de cada zona.\citep{luyben2014} O sensoriamento é realizado
integralmente por termopares tipo~K, distribuídos conforme a
Tabela~\ref{tab:instrumentacao}.

\begin{itemize}
  \item Forno pré-aquecedor (2~TC): monitoram a zona de pré-aquecimento e
    fornecem o sinal de realimentação à malha única.
  \item Camisa externa dos fornos reatores (3~TC por forno, 6~total):
    medem o perfil axial de temperatura da parede e constituem a variável de
    processo (PV) das malhas PID internas (\emph{inner loops}) de cada zona;
    a variável manipulada (MV) correspondente é a potência fornecida ao
    elemento resistivo daquela zona de aquecimento.
  \item Interior dos reatores catalíticos (2~TC por reator, 4~total):
    instalados diretamente no leito, medem a temperatura de reação e definem
    o setpoint objetivo ($SP_{\mathrm{d}}$) do sistema de controle. A escolha
    da temperatura da camisa como PV das malhas internas, em vez da temperatura
    do leito, decorre de um requisito de segurança operacional: ao regular a
    camisa, o controlador impõe um limite superior à temperatura do forno,
    protegendo a integridade física da planta contra sobreaquecimento. O
    objetivo final de manter a temperatura do leito catalítico no valor
    desejado é perseguido indiretamente, por meio do mecanismo de setpoint
    dinâmico proporcional descrito na Seção~\ref{sec:estrategia}.
\end{itemize}

\begin{table}[ht]
\begin{center}
\caption{Instrumentação de temperatura e malhas de controle da planta piloto.}
\label{tab:instrumentacao}
\begin{tabular}{lcc}
\hline
Unidade & TC tipo~K & Malhas \\
\hline
Forno pré-aquecedor              & 2 & 1 \\
Forno reator 1 (camisa ext.)     & 3 & 3 \\
Forno reator 1 (leito int.)      & 2 & 0 \\
Forno reator 2 (camisa ext.)     & 3 & 3 \\
Forno reator 2 (leito int.)      & 2 & 0 \\
\hline
Total & 12 & 7 \\
\hline
\end{tabular}
\end{center}
\end{table}

% -----------------------------------------------------------------------
\subsection{Sinais de controle das malhas}
\label{sec:sinais_controle}
% -----------------------------------------------------------------------

Esta subseção descreve formalmente as variáveis de processo, variáveis
manipuladas e sinais de referência de cada malha de controle, estabelecendo
a nomenclatura utilizada ao longo do artigo.

\subsubsection{Definição das variáveis por malha}

Cada uma das sete malhas automáticas é caracterizada pela tripla
\(\{SP,\,PV,\,MV\}\). Para as malhas dos fornos reatores, os sinais assumem
os papéis descritos a seguir, com índices $i \in \{1,2\}$ para o reator e
$j \in \{\text{alta, med, baixa}\}$ para a zona axial:

\begin{itemize}
  \item \textbf{Variável de processo} (\textit{process variable}, PV):
    temperatura medida pelo termopar da camisa externa da zona $j$ do
    reator $i$, denotada $T^{(i)}_{\text{cam},j}$.
  \item \textbf{Variável manipulada} (\textit{manipulated variable}, MV):
    potência elétrica aplicada ao elemento resistivo da zona $j$ do reator $i$,
    denotada $P^{(i)}_j$, expressa em percentual da potência nominal (\%).
  \item \textbf{Referência interna} (\textit{setpoint} aplicado ao PID),
    $SP^{(i)}_{\mathrm{ap},j}$: calculada dinamicamente pela lógica de nível
    externo conforme as equações~(\ref{eq:sp_alto_med})--(\ref{eq:sp_baixo}),
    a partir do setpoint desejado $SP_{\mathrm{d}}$ e da temperatura do leito
    $T^{(i)}_{\mathrm{int},j}$.
\end{itemize}

Para a malha do forno pré-aquecedor, o sinal de PV é a temperatura da zona
de pré-aquecimento e a MV é a potência do resistor correspondente; o setpoint
é fixado diretamente pelo operador.

% --- ESPAÇO PARA DISCUSSÃO: Sinais de controle medidos ---
% TODO: Inserir aqui a análise dos sinais registrados durante os ensaios.
% Sugestões de conteúdo:
%   - Faixa de variação observada de cada MV (P^(i)_j) ao longo dos ensaios
%   - Tempos de resposta entre variação de MV e resposta de PV
%   - Comparação entre SP_ap calculado e T_cam medida (seguimento de referência)
%   - Discussão sobre saturação ou limitação de atuador observada nos ensaios
% -----------------------------------------------------------------

\subsubsection{Sinal de controle da potência de aquecimento}
\label{sec:sinal_mv}

A variável manipulada $P^{(i)}_j$ é gerada pela saída do controlador PID
correspondente e enviada ao CLP S7-1200, que a converte em um sinal de
acionamento do elemento resistivo por modulação de fase ou ciclo completo.
O sinal de controle é limitado ao intervalo $[0,\,100]\,\%$ da potência
nominal, garantindo operação segura e evitando sobrecarga dos atuadores.

% --- ESPAÇO PARA DISCUSSÃO: Comportamento do sinal MV ---
% TODO: Descrever aqui o comportamento observado do sinal de potência
%       (MV) durante os ensaios de aquecimento. Por exemplo:
%   - O sinal de MV atinge saturação (100%) durante o aquecimento inicial?
%   - Como o MV evolui após atingir o regime estacionário?
%   - Há diferença entre as zonas alta, média e baixa?
%   - Incluir figura: P^(i)_j vs. tempo (sinal de controle ao longo do ensaio)
% -----------------------------------------------------------------

\begin{figure}[ht]
\begin{center}
% === INSERIR FIGURA: Sinais de controle (MV) vs. tempo ===
% \includegraphics[width=8.4cm]{sinais_mv.eps}
\fbox{\parbox{7.5cm}{\centering\vspace{1.5cm}
{\small [FIGURA A INSERIR]\\[4pt]
Sinais de controle $P^{(i)}_j$ (MV) vs.\ tempo para as três zonas\\
do reator durante ensaio representativo}\vspace{1.5cm}}}
\caption{Sinais de controle (potência de aquecimento, MV) das três zonas
do forno reator ao longo de um ensaio representativo.}
\label{fig:sinais_mv}
\end{center}
\end{figure}

\subsubsection{Sinal de erro e atuação do controlador PID}
\label{sec:sinal_erro}

O erro de controle da malha interna é definido como

\begin{equation}
  e^{(i)}_j(t) = SP^{(i)}_{\mathrm{ap},j}(t) - T^{(i)}_{\mathrm{cam},j}(t),
  \label{eq:erro_malha}
\end{equation}

onde $SP^{(i)}_{\mathrm{ap},j}(t)$ é a referência calculada pelo nível externo
e $T^{(i)}_{\mathrm{cam},j}(t)$ é a temperatura da camisa medida naquele instante.
O controlador PID discreto implementado no CLP opera na forma padrão com
anti-windup, gerando a ação de controle

% Equação quebrada em multline para caber na coluna de 8,4 cm (SBA/IFAC)
\begin{multline}
  P^{(i)}_j(k) = K_p\,e^{(i)}_j(k)
    + \frac{K_p}{T_i}\sum_{n=0}^{k}e^{(i)}_j(n)\,\Delta t \\
    + K_p\,T_d\,\frac{e^{(i)}_j(k)-e^{(i)}_j(k-1)}{\Delta t},
  \label{eq:pid_discreto}
\end{multline}

onde $K_p$, $T_i$ e $T_d$ são os parâmetros de ganho proporcional, tempo
integral e tempo derivativo, respectivamente, e $\Delta t$ é o período de
amostragem do CLP.

% --- ESPAÇO PARA DISCUSSÃO: Parâmetros PID e sintonização ---
% TODO: Inserir aqui os valores de K_p, T_i, T_d utilizados na sintonia final,
%       e descrever o método de sintonia adotado (ex.: Ziegler-Nichols,
%       sintonia empírica, etc.). Discutir se a mesma sintonia foi aplicada
%       a todas as zonas ou se foi necessário ajuste individual por zona.
% -----------------------------------------------------------------

% -----------------------------------------------------------------------
\subsection{Estratégia de controle adotada}
\label{sec:estrategia}
% -----------------------------------------------------------------------

\subsubsection{Comportamento térmico observado e limitação do PID simples}

A primeira abordagem investigada consistiu em uma malha simples PID por zona,
na qual o controlador atuava diretamente sobre a potência de aquecimento do
forno tendo como variável de processo a leitura dos termopares das camisas
dos reatores. Essa estrutura revelou-se insatisfatória em operação real
devido a um comportamento térmico assimétrico: a relação entre a temperatura
da camisa e a temperatura efetiva no interior do reator não é fixa, variando
de forma significativa conforme o estado reacional do sistema.\citep{luyben2014}

O fenômeno pode ser descrito da seguinte forma. Com a reação em curso, o
caráter fortemente endotérmico do processo mantém o leito interno
significativamente mais frio que a camisa externa. Já na ausência de reação,
como durante o aquecimento inicial ou após a interrupção da alimentação,
o leito pode atingir temperatura superior à da camisa, pois não há consumo
endotérmico para dissipar o calor acumulado. Essa inversão do sentido do
gradiente térmico em função do estado reacional inviabiliza uma sintonia PID
única válida em toda a faixa de operação.\citep{luyben2014,essien2019}
O comportamento comparativo nos dois regimes é ilustrado nas
Figuras~\ref{fig:pid_com_reacao} e~\ref{fig:pid_sem_reacao}.

\begin{figure}[ht]
\begin{center}
% === INSERIR FIGURA: PID simples COM reação ===
% \includegraphics[width=8.4cm]{pid_com_reacao.eps}
\fbox{\parbox{7.5cm}{\centering\vspace{1.5cm}
{\small [FIGURA A INSERIR]\\[4pt]
PID simples, com reação em curso:\\temperaturas da camisa e do leito vs.\ tempo}\vspace{1.5cm}}}
\caption{PID simples com reação em curso: gradiente entre camisa e leito
decorrente da endotermia e afastamento do leito em relação a $SP_{\mathrm{d}}$.}
\label{fig:pid_com_reacao}
\end{center}
\end{figure}

\begin{figure}[ht]
\begin{center}
% === INSERIR FIGURA: PID simples SEM reação ===
% \includegraphics[width=8.4cm]{pid_sem_reacao.eps}
\fbox{\parbox{7.5cm}{\centering\vspace{1.5cm}
{\small [FIGURA A INSERIR]\\[4pt]
PID simples, sem reação (aquecimento a vazio):\\inversão do gradiente e sobreaquecimento}\vspace{1.5cm}}}
\caption{PID simples sem reação: inversão do gradiente térmico e sobreaquecimento
do leito em relação a $SP_{\mathrm{d}}$.}
\label{fig:pid_sem_reacao}
\end{center}
\end{figure}

% --- ESPAÇO PARA DISCUSSÃO: Sinais de controle no PID simples ---
% TODO: Discutir o comportamento dos sinais de controle (MV e PV) observados
%       nos ensaios com PID simples, em ambos os regimes (com e sem reação).
%       Relacionar o comportamento da MV (potência) com o comportamento
%       da PV (temperatura de camisa) e do leito interno.
% -----------------------------------------------------------------

\subsubsection{Controle por setpoint dinâmico proporcional}

Para superar as limitações identificadas, adotou-se uma estrutura de controle
em dois níveis. No nível interno, um controlador PID independente atua em cada
uma das três zonas de aquecimento de cada forno reator, tendo como variável de
processo a temperatura da camisa externa correspondente e como variável
manipulada a potência do elemento resistivo daquela zona. No nível externo, uma
lógica de \emph{setpoint dinâmico proporcional} calcula, a cada ciclo de
controle, o valor de referência a ser enviado a cada PID interno, compensando
automaticamente o gradiente térmico entre a camisa e o interior do reator. A
arquitetura completa é ilustrada na Figura~\ref{fig:blocos_controle}.

\begin{figure}[ht]
\begin{center}
% === INSERIR FIGURA: Diagrama de blocos do novo controlador ===
% \includegraphics[width=8.4cm]{diagrama_blocos_controle.eps}
\fbox{\parbox{7.5cm}{\centering\vspace{1.5cm}
{\small [FIGURA A INSERIR]\\[4pt]
Diagrama de blocos, setpoint dinâmico proporcional:\\nível externo: $SP_{\mathrm{ap}}$ via (3)/(4);\\nível interno: PID por zona de camisa}\vspace{1.5cm}}}
\caption{Diagrama de blocos do controlador por setpoint dinâmico proporcional:
nível externo calcula $SP_{\mathrm{ap}}$ via (3) e (4); nível interno executa
os PIDs das três zonas da camisa.}
\label{fig:blocos_controle}
\end{center}
\end{figure}

O princípio da estratégia é o seguinte. Em vez de enviar diretamente ao PID
interno o setpoint de temperatura desejado no leito, $SP_{\mathrm{d}}$,
calcula-se uma correção proporcional ao desvio entre $SP_{\mathrm{d}}$ e a
temperatura medida no interior do reator, $T_{\mathrm{int}}$. O setpoint
aplicado ao controlador PID da camisa, $SP_{\mathrm{ap}}$, é dado por

\begin{equation}
  SP_{\mathrm{ap}} = SP_{\mathrm{d}} + \left(SP_{\mathrm{d}} - T_{\mathrm{int}}\right),
  \label{eq:sp_dinamico}
\end{equation}

em que $SP_{\mathrm{d}}$ é a temperatura desejada no interior do reator
(\(^\circ\)C) e $T_{\mathrm{int}}$ é a leitura do termopar interno
correspondente (\(^\circ\)C). A equação~(\ref{eq:sp_dinamico}) pode ser
reescrita de forma equivalente como

\begin{equation}
  SP_{\mathrm{ap}} = 2\,SP_{\mathrm{d}} - T_{\mathrm{int}},
  \label{eq:sp_simplificada}
\end{equation}

mostrando que o setpoint aplicado à camisa é calculado simetricamente em
relação ao valor desejado: quando o leito está mais frio que $SP_{\mathrm{d}}$
(situação típica durante a reação), $SP_{\mathrm{ap}}$ sobe acima de
$SP_{\mathrm{d}}$, aumentando o fluxo de calor; quando o leito está mais
quente (situação típica sem reação), $SP_{\mathrm{ap}}$ é reduzido,
evitando o sobreaquecimento.

% --- ESPAÇO PARA DISCUSSÃO: Dinâmica do sinal SP_ap ---
% TODO: Descrever a dinâmica observada do sinal SP_ap calculado ao longo
%       dos ensaios. Questões a abordar:
%   - Qual a faixa de variação observada de SP_ap em relação a SP_d?
%   - O sinal SP_ap é suave ou apresenta oscilações? Por quê?
%   - Há saturação de SP_ap acima do limite de temperatura do forno?
%   - Incluir figura: SP_ap, T_cam e T_int vs. tempo (três curvas sobrepostas)
%     para ilustrar o mecanismo de compensação dinâmica.
% -----------------------------------------------------------------

\subsubsection{Distribuição axial dos termopares internos e das malhas}

Cada reator catalítico possui dois termopares internos ao leito: um na região
superior, $T_{\mathrm{alto}}$, e um na região inferior, $T_{\mathrm{baixo}}$.
As malhas das zonas alta e média da camisa utilizam $T_{\mathrm{alto}}$ como
$T_{\mathrm{int}}$ na equação~(\ref{eq:sp_simplificada}), resultando em

\begin{equation}
  SP_{\mathrm{ap,\,alta}} = SP_{\mathrm{ap,\,med}} =
  2\,SP_{\mathrm{d}} - T_{\mathrm{alto}},
  \label{eq:sp_alto_med}
\end{equation}

enquanto a malha da zona baixa utiliza exclusivamente $T_{\mathrm{baixo}}$,

\begin{equation}
  SP_{\mathrm{ap,\,baixa}} = 2\,SP_{\mathrm{d}} - T_{\mathrm{baixo}},
  \label{eq:sp_baixo}
\end{equation}

garantindo que cada zona da camisa seja compensada pelo gradiente térmico real
da região axial do leito que ela aquece.\citep{luyben2014}

% --- ESPAÇO PARA DISCUSSÃO: Perfil axial dos sinais de controle ---
% TODO: Discutir as diferenças entre os sinais das zonas alta, média e baixa.
%   - Os valores de SP_ap diferem significativamente entre as zonas?
%   - O perfil axial de temperatura (T_alto vs. T_baixo) é uniforme?
%   - A estratégia de associar T_alto às zonas alta+média e T_baixo à
%     zona baixa se mostrou adequada? Há necessidade de refinamento?
% -----------------------------------------------------------------
